

\chapter{Results and Discussion}

In this chapter, we show how well the proposed model works for identifying various depression severity levels in tweets. 

We used Python version 3 for the experimental setup. It helped us easily use different deep learning tools and libraries. To understand both word and sentence-level features, the model uses advanced methods like hierarchical attention and BERT-based embeddings. Table \ref{tab:hyper_parameters} shows the hyperparameter values we used to fine-tune and improve the performance of model.

\begin{table}[htbp]
	\centering
	\caption{Hyperparameters}  % Removed space between Hyper and parameters
	\label{tab:hyper_parameters}
	\small
	\begin{tabular}{
			>{\RaggedRight}p{6.0cm}  % Fixed column specification
			>{\RaggedRight}p{6.0cm}
		}
		\toprule
		\textbf{Parameters} & \textbf{Values} \\
		\midrule
		Optimizer & Adam \\
		\hline
		\addlinespace
		Loss & Sparse Categorical Cross-Entropy \\  % Fixed hyphenation
		\hline
		\addlinespace
		Epochs & 10 \\
		\hline
		\addlinespace
		Activation & Softmax (output layer) \\
		\hline
		\addlinespace
		Learning rate & 0.00005 \\
		\hline
		\addlinespace
		Batch size & 32 \\
		\bottomrule  % Added missing bottom rule
	\end{tabular}
\end{table}

\section{Performance metrics} 

In this section, we explain the performance metrics used to check the performance of model. These metrices include accuracy, precision, recall, and F1-score.